\honest{My Kind of Reflection}{Eliezer Yudkowsky}{2008}

In ``Where Recursive Justification Hits Bottom'' I concluded that it is fine to use induction to judge induction, and Occam's razor to judge Occam's razor. I am ``far from the first person'' to think about this. A commenter, Chris Hibbert, compared my view to Bartley's pancritical rationalism, so here are the features of my view, which ``may or may not'' be shared by other philosophers. \nb{Hibbert's comment also said where the views part, preferring ``the emphasis on process'' to Yudkowsky's ``focus on justification.'' Bartley dropped the demand that beliefs be justified. The previous post keeps it: ``Everything, without exception, needs justification.''}

First, all of this comes from trying to design a self-modifying AI that applies its reasoning principles to itself while rewriting its code; using induction to license induction really means an inductive AI considering a rewrite of its own induction. If you would not want the AI to stop using induction, your philosophy ``had better not label induction as unjustifiable.'' \nb{The test rules out the classic position. Hume held that inferences from experience come from custom, not reasoning, and that only ``a fool or madman'' would reject experience as a guide.}

Second, the true Way generally turns out to be naturalistic: the AI should treat its own transistors like transistors in the environment, not as a special case, which is why I called such questions ordinary. I strongly suspect that a correctly built AI will not need to do anything special when it changes its own Occamian reasoning. \nb{Quine proposed in 1969 that the theory of knowledge be studied as part of natural science, and answered the charge of circularity.}

Third, reflective consistency is not a goal in itself. Five accurate maps of a city agree, and so do five copies of one fantasy map. The goal is to win; consistency comes along, as the parts of a map someone keeps improving become consistent. The AI is not trying to be reflectively consistent but to optimize the expected utility of its source code, using its current mind's anticipations. Fourth, Jaynes said to use all the information you have, so use induction when you think about induction. \nb{Both points were in the previous post.}

Fifth, take an aggressive posture toward the truth, not a defensive one toward a questioning philosopher: inspect your foundations to improve them or because you fear they are wrong, not to look fair; I deprecate mere dutiful doubts. We did improve one: intuition finds Thor simpler than Maxwell's equations, and information theory ranks the equations simpler. Sixth, the improvements should add up to normality: 2 + 2 is still 4, not something exciting like ``fish.''

Seventh, ``Does induction work?'' is a question an AI may answer by induction. ``Why does induction work?'' is still a mystery, where strange speculation may be needed for now. But claiming that the universe is not regular, that induction does not work, is 2 + 2 = 3 territory, trying to be interesting instead of correct. And answering ``why'' with ``Because induction works'' is circular, while ``Because I believe induction works'' is magical thinking. \nb{Hume's charge of circularity was made against the first question, and answering it by induction is the rule-circular reply the previous post defends.}

Last, a loop of justification through the meta-level is not circular logic, which is a forbidden thing on the object level, while forbidding reflective coherence does not sound wise; but I have not yet formalized the difference; my theory is ``something I'm trying to work out, not something I have in hand.'' \nb{Philosophers call the two kinds rule-circular and premise-circular arguments, and some, such as Van Cleve and Papineau, had argued that the first need not be vicious. In 2013 Yudkowsky and Marcello Herreshoff tried to formalize a self-modifying agent and met a ``L\"obian obstacle,'' which they could get around only by technical methods.}
